\documentclass[11pt]{article}
\usepackage[a4paper,margin=25mm]{geometry}
\usepackage{amsmath}
\usepackage{booktabs}
\usepackage{hyperref}

\title{Academic Workstation Synthetic LaTeX Fixture}
\author{Public acceptance test}
\date{2026}

\begin{document}
\maketitle

\begin{abstract}
This document is synthetic test data. It exercises compilation, references, citations,
figures, tables, equations, and PDF page rendering; it is not a research result.
\end{abstract}

\section{Purpose}\label{sec:purpose}
The acceptance route keeps the editable \LaTeX{} source and treats the PDF as a derivative QA
surface. The fixture is intentionally small, deterministic, and safe to publish. The workflow
described here follows the evidence pattern in the Academic Workstation matrix
\cite{syntheticAcceptance}.

\section{Signal fixture}\label{sec:signal}
Figure~\ref{fig:path} shows the input figure drawn from a text source rather than a binary asset.
The synthetic signal is represented by
\begin{equation}
  y(t) = a\sin(2\pi f t + \phi) + \varepsilon(t),
  \label{eq:signal}
\end{equation}
where $a$ is amplitude, $f$ is frequency, $\phi$ is phase, and $\varepsilon$ is a synthetic
perturbation. Equation~\ref{eq:signal} is deliberately labeled so cross-reference resolution
is visible in the compile log.

\begin{figure}[ht]
  \centering
  \setlength{\unitlength}{1mm}
\begin{picture}(105,42)
  \linethickness{0.4pt}
  \put(8,8){\line(1,0){88}}
  \put(8,8){\line(0,1){27}}
  \put(8,35){\line(1,0){88}}
  \put(96,8){\line(0,1){27}}
  \put(8,8){\line(1,1){18}}
  \put(26,26){\line(1,-1){18}}
  \put(44,8){\line(1,1){18}}
  \put(62,26){\line(1,-1){18}}
  \put(80,8){\line(1,1){16}}
  \put(45,1){\makebox(0,0){synthetic signal path}}
\end{picture}

  \caption{A synthetic signal path drawn with the LaTeX picture environment.}
  \label{fig:path}
\end{figure}

\clearpage
\section{Tabular fixture}\label{sec:table}
Table~\ref{tab:measurements} contains deterministic values used only to exercise table
generation. The values are not measurements and must not be interpreted as experimental data.

\begin{table}[ht]
  \centering
  \caption{Synthetic measurements for acceptance testing.}
  \label{tab:measurements}
  \begin{tabular}{@{}lrrr@{}}
    \toprule
    Group & Count & Score & Flag \\
    \midrule
    Alpha & 4  & 51.4 & PASS \\
    Beta  & 6  & 52.1 & PASS \\
    Gamma & 8  & 52.8 & PASS \\
    Delta & 10 & 53.5 & PASS \\
    \bottomrule
  \end{tabular}
\end{table}

The table is cross-referenced from Section~\ref{sec:table}, while the figure and equation
remain available from Section~\ref{sec:signal}. Keeping these links live distinguishes a
structured source from a flattened page image.

\clearpage
\section{Acceptance notes}\label{sec:notes}
The compiler records its selected engine, command diagnostics, output PDF, and SHA-256 hashes.
The PDF QA stage then checks page count and geometry, renders each page, and requires one visual
observation per page. A missing visual observation is a failed gate, not an implicit pass.

\begin{itemize}
  \item Source: editable \texttt{main.tex}, \texttt{figure.tex}, and \texttt{references.bib}.
  \item Structural checks: equation, table, figure, citation, and cross-reference.
  \item Derivative checks: PDF structure, page rendering, and page-by-page visual review.
\end{itemize}

\clearpage
\section{Reproducibility boundary}\label{sec:boundary}
The public fixture does not include private workstation paths, generated binaries, credentials,
or user documents. The exact TeX distribution and PDF tools are recorded in private evidence at
runtime. This separation lets the repository demonstrate the acceptance contract without
claiming that every host has the same native capability.

\bibliographystyle{plain}
\bibliography{references}

\end{document}
